% SGA 3，第 VII_B 篇，英文重建稿——第 1 轮。
% 范围：权威底本分篇第 93--94 页，第 5.2.4 节全文，止于注 5.2.4.A 之前。
% 合并阅读本页码： 606--607.
% 再版印刷页码： 595--596.
% 本范围内原版页边定位： 561.
% 本范围内编者注： 160.
% 权威底本： Polo--Gille Exp7B-13oct24.pdf,
% SHA-256 D303690E0195654BBE7C3DAA5D89EA23E24CD709A3575217234AC19AE8836694.
% 曾参考 Jacob Reinhold 的 CC BY 4.0 Markdown 作为比对谱系。
% 原英文 Loop 1 阶段说明：各行正合的图在 Loop 1 中以直接取自 PDF 的精确 PNG 图像保留，原计划在 Loop 2 作原生重建。当前中文工程已采用原生图。

\subsubsection*{5.2.4}
\phantomsection
\label{sga3:VIIB:5.2.4}

对 \(k\) 上的每个伪紧向量空间 \(V\)，用
\[
  {}_{\pi}V
  =
  V\widehat\otimes_{\pi}k
\]
表示由 \(V\) 沿 \(x\mapsto x^p\) 扩张标量所得的空间。\footnote{\label{sga3:VIIB:nde:160}由于 \(k\) 为完美域，可以将 \({}_{\pi}V\) 与阿贝尔群 \(V\) 等同，其上 \(k\) 按 \(\lambda\cdot v=\lambda^{1/p}v\) 作用。}
于是有如下各行正合的交换图：

\begin{figure}[H]
  \centering
  % SGA3-VII-DIAG-135
  % SGA3-VII-DIAG-135
% Authority: Exp7B-13oct24.pdf, local p. 93.
% Crop: figures/Exp7B_localp093_pi_exact_rows_diagram.png
% Locator: tex/components/89_expose_VIIB_section524_frobenius_twist_induction_en.tex, Loop 1 line 32.
\begin{tikzcd}[row sep=2.8em,column sep=2.6em]
0 \arrow[r]
  & {}_\pi\overline{I^2} \arrow[r,"\alpha"] \arrow[d,"u"']
  & {}_\pi I \arrow[r,"\beta"] \arrow[d,"v"']
  & {}_\pi\omega \arrow[r] \arrow[d,"w"']
  & 0 \\
0 \arrow[r]
  & \overline{J^2} \arrow[r,"\gamma"]
  & J \arrow[r,"\delta"]
  & \overline\omega \arrow[r]
  & 0\, ,
\end{tikzcd}

\end{figure}

\noindent
这里
\[
  \overline\omega=J/J^2,
\]
且映射 \(u,v,w\) 由线性映射
\[
  {}_{\pi}B\longrightarrow B^p,
  \qquad
  x\widehat\otimes a\longmapsto x^pa.
\]
诱导。由于 \(u\) 和 \(v\) 为满射，\(w\) 也是满射，其核为
\[
  {}_{\pi}I_1=\ker(v).
\]
的像 \({}_{\pi}\omega_1\)。

置
\[
  J_n=\{\,x\in J\mid x^{p^n}=0\,\},
  \qquad
  \overline\omega_n=\delta(J_n).
\]
则对每个 \(n\geq0\)，
\[
  J_n=v({}_{\pi}I_{n+1}),
  \qquad
  \overline\omega_n=w({}_{\pi}\omega_{n+1}).
\]

\phantomsection
\label{sga3:VIIB:5.2.4:delta-source-note}
\noindent\textit{底本说明。}\enspace
法文 PDF 印作 \(\overline\omega_n=\delta(I_n)\)，但所列正合行在 \(J\) 上定义了 \(\delta\)，而紧接前面的定义引入了 \(J_n\subset J\)。上文采用类型正确的 \(\delta(J_n)\)；随后所有恒等式都与此读法相符。

因此，由 \SGARefLink{sga3:VIIB:5.2.1:section-proposition}{5.2.1} 的截面 \(\sigma\) 诱导的截面
\[
  {}_{\pi}\sigma:{}_{\pi}\omega\longrightarrow{}_{\pi}I
\]
下降为截面
\[
  \tau:\overline\omega\longrightarrow J
\]
它与 \(J\) 和 \(\overline\omega\) 的滤过相容。由于 \(B^p\) 的高度至多为 \(r-1\)，归纳假设表明此截面诱导同构
\[
  \psi:
  B''
  =
  k[[\overline\omega]]
  \big/
  \bigl((x^{p^n})\bigr)_{x\in\overline\omega_n}
  \xrightarrow{\ \sim\ }
  B^p.
\]
但 \(B''\) 与 \(B'^p\) 等同，\(\psi\) 与 \(\phi\) 诱导的态射 \(B'^p\to B^p\) 等同。因此，定理在 \(B\) 高度有限时得证。
